Quando o sol se pôs, os hóspedes que estavam passeando pela cidade começaram a voltar para o hotel. O leão-marinho estava na recepção, seus longos bigodes caídos. ``Meu cachimbo sumiu!'', ele berrou. ``Foi um presente que ganhei quando me tornei membro honorário do Clube de Bocha do Ártico.'' A coruja grasnou em desespero, segurando uma foto de seu avô. ``Meu monóculo sumiu! É a única coisa que me resta dele!'' A girafa chorou que seu cachecol vermelho havia desaparecido. ``É a única lembrança da minha primeira vez esquiando no Himalaia!'' O coelho pulava para cima e para baixo, procurando seu relógio de pulso. ``Eu preciso desesperadamente do meu relógio! Como vou saber agora quando é a hora do meu remédio?!'' O porco suava de ansiedade. ``Meu jogo de futebol está prestes a começar!'' E bem atrás dele estava um urso, que segurava as calças com as patas. ``Meus suspensórios sumiram!'', ele rosnou. ``Como minhas calças vão ficar no lugar sem eles?''

A girafa, visivelmente aborrecida, acrescentou: ``Precisamos encontrar o ladrão! Eles não podem roubar nossos objetos de valor e sair impunes!'' Sua frustração foi ecoada pelos outros, todos ansiosos para resolver o crime.

O coelho, aproveitando a determinação de todos, apontou: ``Quando desci para a farmácia, vi o lobo perambulando pelas escadas e corredores.'' Foi então que o lobo apareceu e, com uma tosse engasgada, disse em voz alta: ``Cof! Ei! Eu também sou uma vítima aqui! Estava perambulando por aí procurando meu vale-spa, que ganhei na recepção!''

Henrietta interveio rapidamente, pedindo a todos que se acalmassem. ``O Sr. Lobo foi de fato premiado com um vale-spa na sua chegada.''

Então, a avestruz apareceu, segurando um bilhete em seu bico. ``Acho que encontrei seu vale, Sr. Lobo.''

\begin{center}
\includegraphics[width=0.2\textwidth]{images/ostrich.png}
\end{center}

O coelho, com um olhar astuto, apontou: ``Podemos então concluir que o lobo não foi uma vítima dos desaparecimentos...''

Os animais começaram a se voltar uns contra os outros. A coruja acusou o coelho de ser muito esperto e inquieto. Além disso, ela acrescentou, ``Ele está apontando o dedo para o lobo para tirar o foco de si mesmo.'' O coelho retrucou, ``Isso é uma acusação leviana!'' Nesse momento, um galo de crista vermelha, aproveitou a confusão e cantou em voz alta: ``Não confiem nela!'' Ele apontou a asa para a raposa, acusando-a de ser muito astuta e sorrateira. A raposa, por sua vez, apontou a pata para o porco, alegando que ele era desajeitado e bagunceiro. O porco, frustrado, apenas bufou em descrença.

Quando a confusão atingiu seu auge, uma garota chamada Anne chegou com seu Border collie peludo, Scotch. Ela ouviu atentamente as histórias dos animais, enquanto Scotch, por outro lado, estava cheirando tudo.

``Podemos ajudar'', disse Anne com confiança. ``Scotch pode seguir o rastro dos seus objetos perdidos.''

A raposa e o lobo trocaram um olhar astuto. ``Por que vocês não começam pelo quarto do coelho?'', sugeriu a raposa. ``Ele parece terrivelmente nervoso.''

``Podem ir!'' o coelho exclamou, estufando o peito. ``Podem ver meu quarto. Eu não tenho nada a temer, e vocês vão se decepcionar!''

Todos foram para o quarto do coelho. Scotch começou a cheirar tudo: a porta, o chão, o guarda-roupa... e de repente, eles viram um monóculo brilhante caído no chão, ao lado do elevador de carga. O lobo soltou um bufo satisfeito. ``Bem, parece que encontramos o que estávamos procurando.''

\clearpage
\includepdf[pages={1},
            pagecommand={\thispagestyle{fancy}}, % Apply your custom 'fancy' style
            fitpaper=true,
            noautoscale
           ]{rabbitpanic.pdf}


O coelho gaguejou, ``Eu não sei como isso veio parar aqui!'' Ele suou frio, ficou pálido, e parecia que ia desmaiar.

Mas Scotch latiu mais uma vez. Ele se ergueu sobre as patas traseiras, com as patas dianteiras apoiadas na parede para se equilibrar. Com o nariz levantado, ele cheirou o elevador avidamente, seu corpo tenso de curiosidade e excitação.

``Parece que o rastro leva a este pequeno elevador'', Anne acrescentou. ``Para onde ele vai?''

Bernard respondeu prontamente: ``Ele vai para todos os quartos do hotel.''

Henrietta acrescentou: ``Mas há um número infinito de quartos... levará uma eternidade para encontrar os itens perdidos!''

Scotch latiu novamente e começou a puxar Anne para fora do quarto. Ele a conduziu pelo corredor até um carrinho de serviço de quarto que estava parado no meio. Um pano branco cobria os itens na prateleira inferior, mas o nariz de Scotch estava trabalhando a todo vapor. Com o cão curioso cheirando com afinco, Anne levantou o pano com cuidado para ver o que havia por baixo.

Ela ficou surpresa ao encontrar um pote cheio de biscoitos de cachorro. Ela riu, ``Haha! Scotch, seu guloso!''

Nesse momento, Bernard se aproximou, abriu o pote e lhe deu três biscoitos. Scotch latiu, como se estivesse agradecendo pelo agrado.

%\begin{center}
%\includegraphics[width=0.55\textwidth]{images/biscuits.png}
%\end{center}

%Após seu merecido lanche, Scotch então disparou de volta para o saguão, com todo o grupo seguindo logo atrás dele. Ele encontrou outra porta para o elevador de carga, sentou e latiu. 
Após seu merecido lanche, Scotch conduziu o grupo. Procuraram quarto após quarto, mas eles não conseguiram encontrar uma única pista. À medida que a esperança dos animais começava a se esvair, a pura impossibilidade da tarefa -- encontrar os itens perdidos nessa infinidade de quartos -- se tornava evidente. Quando eles estavam prestes a voltar para o saguão, Scotch encontrou outra porta para o elevador de carga, sentou e latiu.
Anne olhou da porta para o grupo de animais, com um olhar de quem descobriu a resposta. ``Acho que entendi agora'', ela disse. ``Esta porta leva ao mesmo elevador, não é?''

Bernard imediatamente acrescentou: ``Exato! Para mover os pertences dos hóspedes, precisávamos usar esta entrada para o elevador para que pudéssemos escolher para qual quarto enviá-los.''

\clearpage
\includepdf[pages={1},
            pagecommand={\thispagestyle{fancy}}, % Apply your custom 'fancy' style
            fitpaper=true,
            noautoscale
           ]{biscuit.pdf}

\clearpage
\includepdf[pages={1},
            pagecommand={\thispagestyle{fancy}}, % Apply your custom 'fancy' style
            fitpaper=true,
            noautoscale
           ]{manual.pdf}

Anne apontou, pensativa: ``Está começando a ficar claro. Só precisamos descobrir o funcionamento intrincado do sistema do elevador.''

Bernard perguntou: ``Temos um mapa do sistema do elevador?''

Henrietta explicou que eles só tinham um manual, e Anne comentou: ``Não poderia existir um mapa assim, ou ele seria infinitamente grande!'' Bernard quase riu e disse: ``Vou procurar o manual. Já volto!''

Bernard rapidamente voltou com o manual nas mãos. Ele o entregou a Anne, que imediatamente começou a folhear suas páginas. ``Ah! Achei!'' ela exclamou com entusiasmo. ``Os engenheiros usaram a Função de Emparelhamento de Cantor para mapear o número do quarto e o andar com a posição do elevador em sua trilha infinita.'' 
Ela mostrou a Henrietta uma fórmula no livro, explicando como ela realizava o mapeamento. %mapeava o número do quarto e o andar para um único número que dizia ao elevador para onde ir. 
Ela também alertou: ``Temos que lembrar que o índice aqui começa com zero!''
%Ela mostrou a Henrietta uma fórmula no livro: $n=(i+j)(i+j+1)/2+j$. Anne explicou: ``$n$ é a posição do elevador, $i$ é o número do quarto e $j$ é o andar.'' Ela também alertou: ``Temos que lembrar que o índice aqui começa com zero!''

Todos olharam para Anne com alívio, apesar de não conseguirem entender a matemática que ela usou. Ela abriu a bolsa e pegou uma calculadora científica. Depois de algumas teclas, exclamou: ``Entendi! Vamos, Scotch! Vamos encontrar os tesouros desaparecidos!''
